% from main.tex
%%%%%%%%%%  scan idx 549  |  folio 536  %%%%%%%%%%
Poincaré $\chi(F)\in\mathbb Z$ de $F$ sur $X$. De 2.2.1 on conclut immédiatement :

\textbf{Proposition 2.3.1.} Pour tout schéma propre $X$ sur $k$, et pour tout $y\in\Filt_j(X)$, l'application
\[
S_y:\Pic(X)\longrightarrow\mathbb Z\hookrightarrow\mathbb Q
\]
définie par
\[
\cl_P(L)\rightsquigarrow g_*(\cl^\circ(L)y)=\chi(\cl^\circ(L)y)
\]
est une application polynôme de degré $\leq j$.

Lorsque $L_1,\ldots,L_n$ sont des $\mathcal O_X$-modules inversibles et $F$ un $\mathcal O_X$-module cohérent tel que $\dim(\Supp(F))=j$, alors 2.3.1. implique en particulier que la caractéristique d'Euler-Poincaré
\begin{equation}\tag{2.3.2}
\chi(L_1^{\otimes m_1}\otimes\cdots\otimes L_n^{\otimes m_n}\otimes F)
\end{equation}
est un polynôme numérique dans les variables $m_1,\ldots,m_n\in\mathbb Z$, de degré $\leq\dimdot\operatorname{supp.}F$, parfois appellé polynôme de Snapper pour $F$ et la suite $L_1,\ldots,L_n$.

\subsection*{2.4}
Rappelons [2] que le foncteur de Picard $\Pic_{X/k}$ de $X$ sur $k$ est le faisceau en groupes dans la topologie ($\fpqc$) de $k$ associé au préfaisceau
\[
T\rightsquigarrow \Pic(X\times_kT)\qquad (T\in\operatorname{Ob}(\Sch/k)).
\]
Lorsque $k$ est algébriquement clos, on sait que l'homomorphisme canonique
\[
\Pic(X)\longrightarrow\Pic_{X/k}(\Spec(k))
\]
est un isomorphisme.



% Local continuation macros
\providecommand{\Filt}{\operatorname{Fil}}
\providecommand{\fpqc}{\mathrm{fpqc}}
\providecommand{\Sch}{\operatorname{Sch}}
\renewcommand{\Ob}{\operatorname{ob}}
\providecommand{\Gm}{\mathbb G_m}
\providecommand{\Der}{\operatorname{Der}}
\providecommand{\SL}{\operatorname{SL}}
\providecommand{\Aut}{\operatorname{Aut}}
\providecommand{\Mod}{\operatorname{Mod}}
\providecommand{\Tors}{\operatorname{Tors}}
\providecommand{\Mat}{\operatorname{Mat}}
\providecommand{\len}{\operatorname{long}}
\renewcommand{\id}{\operatorname{id}}
\providecommand{\Spec}{\operatorname{Spec}}
\providecommand{\Picf}{\operatorname{Pic}}
\providecommand{\Ens}{\operatorname{Ens}}
\providecommand{\clp}{\operatorname{cl}_{P}}
\providecommand{\clo}{\operatorname{cl}^{\circ}}
\providecommand{\clb}{\operatorname{cl}_{\bullet}}
\providecommand{\Grb}{\operatorname{Gr}_{\bullet}}
\providecommand{\Gro}{\operatorname{Gr}^{\circ}}
\providecommand{\longO}[1]
